\chapter{感知自适应动态LOD技术}


\section{问题描述与分析}
当前 LOD 技术为平衡渲染质量与性能提供了基础框架，但其在面向动画网格的注视点渲染应用时，面临两大根本性局限，制约了感知驱动渲染效率的进一步提升。

\subsection{问题描述}
在虚拟现实应用中，采用高精度动画网格是提升场景真实感与沉浸感的关键手段，但这同时导致了渲染计算负荷的急剧增加，对维持高帧率与低延迟提出了严峻挑战。
注视点渲染技术通过在人眼不同视觉区域分配差异化渲染资源，为缓解这一矛盾提供了潜在路径。
然而，将现有LOD 技术直接应用于动画网格的注视点渲染时，其效果仍存在显著不足。

首先，基于静态模型的传统 LOD 技术在简化动画网格时，会引入显著的动态视觉失真。
在骨骼动画驱动的高变形区域，简化操作所依据的静态几何误差最小化准则，破坏了模型在时序上的运动连贯性特征，导致简化后的模型在动画播放过程中，出现非自然的几何塌陷或关键帧间的运动跳变。
当用户注视点追踪此类动态目标时，此类失真若落入中央凹视觉区域，将直接损害视觉内容的真实感与沉浸体验。

其次，传统 LOD 的离散切换机制无法实时适配快速变化的视觉注意力，导致资源分配错位：
注视焦点可能出现细节缺失，而外围非关注区域则存在渲染冗余。这种错配既损害了感知质量，也降低了渲染效率。



\subsection{问题分析}

将传统LOD技术直接应用于动画网格的注视点渲染，其适应性存在本质缺陷，现从以下两个方面对该问题进行剖析：

传统LOD的静态简化准则与动画网格的动态视觉保真需求存在根本矛盾。
动画网格的视觉保真度不仅取决于静态几何精度，更关键地依赖于其在时间域上的运动连贯性与特征保持。
现有的 LOD 生成算法本质上是为静态模型设计的空间近似方法，
当应用于动态序列时，其优化目标仅局限于最小化单帧或静态绑定姿态下的几何误差，而忽略了顶点在整个动画周期中所承载的运动显著性。
导致简化后的模型在播放时，于高变形区域产生非自然的几何失真与运动瑕疵。
在注视点追踪动态目标的场景下，此类瑕疵若落入中央凹视觉区域，将显著破坏视觉真实感。
因此，需要提出一种新的LOD技术，引入动态视觉重要性度量，在简化过程中综合考虑几何误差与时序运动特征，以保持动画网格在时间维度上的视觉连贯性。


传统LOD技术在注视点渲染中的适应性存在本质缺陷。
现有方案主要沿袭游戏引擎中的离散LOD切换机制，在简化动画网格时面临双重困境：
在空间维度，离散 LOD 切换会引发几何突变和纹理闪烁；离散的LOD层级无法匹配由注视点引导的、连续变化的感知重要性梯度，导致在注视点附近区域可能出现细节不足，而在外围区域存在渲染冗余；
在时间维度，动态LOD切换难以与快速变化的注视点保持同步，导致视觉焦点区域的细节重建出现可感知的延迟；
此外，传统 LOD 优化以几何误差为评估标准，与注视点渲染所需的心理视觉模型感知优先级存在目标冲突。
从而在注视点附近区域可能出现细节不足，而在外围区域存在渲染冗余，导致质量的可感知损失与渲染资源的浪费。
因此，需要提出一种新的 LOD 技术，利用感知模型衡量不同区域内的模型简化程度，实现更平滑的层次过渡，并平衡视觉感知与计算效率。

综上所述，本文提出一种面向动画网格注视点渲染的感知自适应动态LOD技术，实现在不同注视条件下动态、平滑地调整模型细节以保证视觉感知无损失，同时提升渲染效率。



\section{感知自适应动态LOD技术}

\subsection{总体架构}
针对虚拟现实环境中动画网格渲染面临的动态视觉失真与资源分配错位问题，本章提出一种融合多维层次模型与感知自适应细分LOD的集成化解决方案。该技术以第三章构建的castleCSF-TS感知模型为核心，结合动画网格的时空特性，构建感知自适应的动态层次细节表示体系。

感知自适应动态LOD技术的总体架构如图4.1所示，包含三个协同工作的核心处理模块：

动画感知重要性计算模块：集成castleCSF-TS感知模型与动态运动特征分析，实时生成动画网格的时空重要性图谱。该模块接收原始动画网格、注视点位置和眼动数据，输出每个顶点在时空维度上的感知重要性值。

动态层次构建与优化模块：结合静态网格精简、动态序列优化与感知引导，生成多精度动画表示。该模块采用基于时空感知权重的简化算法，为不同视觉区域生成差异化精度的LOD层次。

注视点驱动的自适应渲染模块：融合离散LOD与连续细分技术，实现基于注视点的平滑细节过渡。该模块根据实时渲染上下文动态调整细节级别，平衡视觉质量与渲染性能。

系统工作流程如下：原始动画网格首先通过结构解耦分离为静态几何与动态动画数据；随后，静态网格精简模块与动态序列优化模块分别对两类数据进行预处理；处理后的数据进入多维层次建模模块，生成多精度层级的统一表示；在实时渲染阶段，注视点跟踪信息驱动感知自适应细分算法，动态调整各区域的细节级别；最终输出与视觉感知一致的高质量渲染帧。


\subsection{时空运动特征驱动的动画网格层次简化算法}

在虚拟现实环境中，动画网格的实时渲染面临着计算复杂度与视觉质量之间的固有矛盾。传统网格简化算法多基于静态几何误差度量进行优化，虽能有效控制单帧几何形变，却常忽略动画序列在时间维度上的运动连贯性。特别在动画网格简化中，基于静态绑定的几何误差最小化准则往往破坏模型在时序上的运动自然性，导致简化后的模型在动画播放过程中出现非物理的几何塌陷或关键帧间的运动跳变。
为解决这一问题本节提出时空运动特征驱动的动画网格层次简化算法，通过构建一个融合动画时序特性与空间结构特征的重要性度量模型，指导生成在动画播放全周期内保持视觉连贯性的多层次细节表示。

（1）时空运动特征重要性度量模型

传统LOD技术在处理动画网格时，主要依赖静态几何误差作为简化准则，忽视动画在时间维度上的视觉重要性变化，导致简化后的模型在播放时，于高变形区域产生非自然的几何失真与运动瑕疵。
为克服这一局限，本节提出时空运动特征重要性度量模型，基于网格在动画序列中表现出的运动规律和结构功能来评估每个顶点的重要性。

对于动画网格中的任意顶点$v$在时间$t$的时空运动特征重要性$I(v,t)$定义为三个关键分量的加权组合：
\begin{equation}
I(v,t) =  \omega_M \cdot M(v,t) + \omega_G \cdot G(v) + \omega_J \cdot J(v)
\end{equation}
其中各分量具体定义如下：

运动显著性分量$M(v,t)$：运动显著性分量量化顶点在整个动画序列中的动态贡献程度，反映了顶点参与重要运动模式的程度，基于顶点在时间窗口内的位移变化率计算：
\begin{equation}
M(v,t) = \frac{1}{T} \sum_{\tau=t-T}^{t} \frac{\| v(\tau) - v(\tau-1) \|}{\Delta t}
\label{eq:motion}
\end{equation}
其中$T$为时间窗口大小，通常取5-10帧。

几何结构分量$G(v)$：几何结构分量评估顶点在维持网格整体形态和表面特性中的重要性，基于顶点曲率$C_v$和法线变化$N_v$计算：
\begin{equation}
G(v) = \alpha \cdot C_v + \beta \cdot \frac{1}{|N(v)|} \sum_{u \in N(v)} (1 - \overrightarrow{n_v} \cdot \overrightarrow{n_u})
\label{eq:geometry}
\end{equation}
其中$N(v)$为顶点$v$的邻接顶点集合，$n_v$和$n_u$分别是顶点$v$和$u$的单位法线向量，$C_v$是顶点$v$的曲率，$\alpha$和$\beta$是权重系数。

关节关联分量$J(v)$：关节关联分量反映了顶点$v$与骨骼关节的接近程度，接近关节的顶点在动画中通常有较大的变形，需要更高的细节保留。
\begin{equation}
J(v) = \exp\left(-\frac{\min_{j \in J} |v - j|^2}{2\sigma^2}\right)
\label{eq:joint_association}
\end{equation}
其中$J$为所有关节的集合。

$ \omega_M, \omega_G, \omega_J$为权重系数，典型配置为 $\omega_M=0.50$, $\omega_G=0.30$, $\omega_J=0.20$，，强调运动特征的主导地位。


（2）时空运动特征驱动的动画网格层次简化算法

基于上述时空运动特征重要性度量模型，本节提出时空运动特征驱动的动画网格层次简化算法，
该算法在传统基于平均偏差的边折叠算法基础上，引入时空运动特征重要性作为加权因子，实现动画网格在时空维度上的自适应简化，在保持视觉关键特征的前提下实现更高压缩率。


在边折叠决策过程中，传统算法主要依赖静态几何误差作为折叠代价的计算基础。设边$e_{ij}=(v_i,v_j)$为待折叠边，其传统的多特征平均偏差$D_{ij}$计算如下：
\begin{equation}
D_{ij} = \frac{1}{3}\left( w_{\text{pos}} \cdot d_{\text{pos}}(v_i, v_j) + w_{\text{norm}} \cdot d_{\text{norm}}(v_i, v_j) + w_{\text{tex}} \cdot d_{\text{tex}}(v_i, v_j) \right)
\label{eq:multi_feature}
\end{equation}
其中$w_{\text{pos}}, w_{\text{norm}}, w_{\text{tex}}$分别为位置、法线和纹理特征的权重系数，$d_{\text{pos}}, d_{\text{norm}}, d_{\text{tex}}$为对应的距离度量函数。位置距离采用欧氏距离，法线距离采用点积运算，纹理距离采用曼哈顿距离。

为将时空运动特征重要性融入简化决策，需要重新定义边折叠的代价函数。对于顶点$v_i$和$v_j$的时空运动特征重要性$I(v_i,t)$和$I(v_j,t)$，首先计算边的平均重要性：
\begin{equation}
I_{\text{avg}}(e_{ij}) = \frac{I(v_i,t) + I(v_j,t)}{2}
\label{eq:avg_importance}
\end{equation}

考虑到重要性值的相对性，对平均重要性进行归一化处理：
\begin{equation}
I_{\text{weight}}(e_{ij}) = \frac{I_{\text{avg}}(e_{ij})}{\max_{v \in V} I(v,t)}
\label{eq:normalized_importance}
\end{equation}
其中$V$为网格顶点集合。归一化后的重要性权重$I_{\text{weight}} \in [0,1]$，反映了该边相对于网格中其他边的重要性程度。

最终，边的折叠代价函数定义为多特征平均偏差与重要性权重的组合形式：
\begin{equation}
\text{Cost}(e_{ij}) = D_{ij} \cdot \left[ 1 + \lambda \cdot (1 - I_{\text{weight}}(e_{ij})) \right]
\label{eq:collapse_cost}
\end{equation}
其中$\lambda \geq 0$为重要性调节系数。当$\lambda=0$时，算法退化为传统的基于多特征平均偏差的简化算法；随着$\lambda$增大，重要性较低的边将获得更小的折叠代价，从而在简化过程中被优先折叠。这种设计确保了高重要性区域得到更好保护，低重要性区域优先简化，实现了简化过程的自适应性。

然而，对于动画网格而言，仅仅考虑静态时刻的重要性度量是不够的。为确保简化后的网格在时间维度上保持运动连贯性，需要引入时序一致性约束。考虑顶点$v$在整个动画序列$T={t_0,t_1,\ldots,t_n}$中的运动轨迹${p(v,t)}$，定义边$e{ij}$折叠后的运动轨迹误差为：
\begin{equation}
E_{\text{temp}}(e_{ij}) = \frac{1}{|T|} \sum_{t \in T} \left| p_{\text{pred}}(v_{\text{new}},t) - p_{\text{actual}}(v_{\text{new}},t) \right|
\label{eq:temporal_error}
\end{equation}
其中$p_{\text{pred}}(v_{\text{new}},t)$为基于原顶点轨迹预测的新顶点位置，通常通过线性插值得到；$p_{\text{actual}}(v_{\text{new}},t)$为实际折叠后产生的新顶点位置。时序误差项$E_{\text{temp}}$的引入确保了简化操作不会破坏动画序列的自然运动规律。

综合静态几何误差和时序运动误差，边的最终折叠代价函数修正为：
\begin{equation}
\text{Cost}_{\text{final}}(e_{ij}) = \text{Cost}(e_{ij}) + \gamma \cdot E_{\text{temp}}(e_{ij})
\label{eq:final_cost}
\end{equation}
其中$\gamma \geq 0$为时序误差权重系数，控制时序一致性在简化决策中的重要性程度。这一复合代价函数同时考虑了空间几何误差、时空运动特征重要性和时间运动连贯性，实现了动画网格简化的多目标优化。

算法的执行过程遵循迭代边折叠的经典范式，但在每个环节都融入了时空运动特征的考量。初始化阶段，首先计算动画网格中每个顶点的时空运动特征重要性$I(v,t)$，提取动画序列中每个顶点的运动轨迹、几何特征和关节关联信，为每条边计算其最终折叠代价，构建边折叠优先队列。

迭代折叠阶段，算法从优先队列中提取当前折叠代价最小的边进行折叠操作。与传统算法不同，本算法在折叠前增加了运动连贯性可行性检验。对于候选边$e_{ij}$，通过分析其在动画序列中的运动轨迹变化，预测折叠操作可能引入的运动失真程度，如果预测的运动失真超过预设阈值$\theta_{\text{motion}}$，即使该边的几何折叠代价很小，也会被拒绝折叠。该策略确保了简化后的网格在动画播放过程中不会出现明显的运动跳变或非自然变形。

边折叠操作执行后，新顶点$v_{\text{new}}$的时空运动特征重要性需要重新计算。新顶点的重要性值不是简单的算术平均，而是通过以下改进的加权平均模型计算：
\begin{equation}
I(v_{\text{new}}, t) = w_i \cdot I(v_i, t) + w_j \cdot I(v_j, t)
\label{eq:weighted_average_importance}
\end{equation}
其中权重系数$w_i$和$w_j$基于几何距离确定，满足$w_i + w_j = 1$。权重的计算考虑了新顶点$v_{\text{new}}$与原始顶点$v_i$、$v_j$之间的空间关系，确保距离较近的原始顶点对新顶点的重要性贡献更大：
\begin{equation}
w_i = \frac{d(v_j, v_{\text{new}})}{d(v_i, v_{\text{new}}) + d(v_j, v_{\text{new}})}, \quad
w_j = \frac{d(v_i, v_{\text{new}})}{d(v_i, v_{\text{new}}) + d(v_j, v_{\text{new}})}
\label{eq:distance_based_weights}
\end{equation}
其中$d(\cdot,\cdot)$表示顶点间的欧氏距离。

边折叠操作改变了网格的拓扑结构，需要相应调整骨骼权重。对于新顶点$v_{\text{new}}$，其骨骼权重向量$w_{\text{new}}$通过原顶点权重进行加权平均得到：
\begin{equation}
w_{\text{new}} = \frac{||w_i|| w_i + ||w_j|| w_j}{||w_i|| + ||w_j||}
\label{eq:bone_weight}
\end{equation}
其中$w_i, w_j$为原顶点$v_i, v_j$的骨骼权重向量，$||\cdot||$表示向量的欧氏范数。基于范数加权的插值方法保持了权重分布的相对比例，避免了简单平均可能导致的失真问题。

对于简化后的网格，需要分析每个骨骼在时间维度上的运动轨迹，识别并剔除运动冗余的关键帧，同时保证重要运动特征的完整保留。对于特定骨骼 \( B_i \) 在关键帧 \( t_j \) 至 \( t_{j+1} \) 时间区间内的运动轨迹，定义其平均位移误差度量函数：：
\begin{equation}
E_{ij} = \frac{1}{N} \sum_{k=1}^{N} I(p_k,t) \cdot \left| \mathbf{T}_{ij}(p_k) - \mathbf{T}_{i(j+1)}(p_k) \right|
\label{eq:motion_weighted_error2}
\end{equation}
其中，\( p_k \) 代表受骨骼影响的顶点集合，\( T_{ij} \) 表示骨骼 \( B_i \) 在关键帧 \( t_j \) 的空间变换矩阵，\( N \) 为该骨骼影响的顶点总数。当$E{ij} < \theta_{\text{key}}$时，判定关键帧$t_{j+1}$在视觉感知上存在冗余，可予以剔除，$\theta_{\text{key}}$ 为位移误差阈值。这种加权误差度量确保在高重要性区域保留更多运动细节，在低重要性区域允许更高的压缩率，实现了运动质量与数据压缩之间的平衡。

为保证压缩后动画序列的连续性和运动平滑性，对保留的关键帧采用Catmull-Rom样条插值算法生成中间过渡帧。插值过程数学表述如下：
\begin{equation}
P(t) = P_1 + \frac{t}{2}\left[-P_0 + P_2 + t\left(2P_0 - 5P_1 + 4P_2 - P_3 + t\left(-P_0 + 3P_1 - 3P_2 + P_3\right)\right)\right]
\label{eq:catmull_rom_interpolation}
\end{equation}
其中$P_0$、$P_1$、$P_2$、$P_3$构成四个连续的控制点序列，$t \in [0,1]$为归一化的时间参数。Catmull-Rom样条具有C1连续性特性，能够生成符合自然运动规律的平滑轨迹，有效维持原始动画的视觉真实感。

纹理数据的优化与几何简化协同进行。基于简化后的网格，重新计算纹理坐标并通过自适应可伸缩纹理压缩算法压缩纹理图像。压缩过程中，对于不同运动特征重要性的区域采用差异化的压缩策略：高重要性区域使用低压缩比保持纹理细节，低重要性区域使用高压缩比减少存储需求。这种分级压缩策略在保证视觉质量的同时显著降低了纹理数据量。

算法迭代执行边折叠操作，达到目标简化率$R_{\text{target}}$后结束。

\subsection{基于感知引导的混合LOD过渡策略}
针对传统LOD技术在动画网格注视点渲染中存在本质性缺陷的问题，本节提出一种基于感知引导的混合LOD过渡策略，将离散简化层次与连续细分技术有机结合，实现动画网格在时空维度上的平滑细节过渡和高效资源分配。该算法框架如图\ref{fig:adaptive_lod}所示。

基于感知引导的混合LOD过渡策略结合了离散 LOD 技术和镶嵌技术。首先，使用离散 LOD 根据注视点位置
初始分阶处理动画网格；其次，在对应 LOD 层次内使用镶嵌技术对注视区域的动画网格
执行自适应镶嵌细分，依据降采样因子动态调整细分程度，使得在中央凹区域的动画网格拥有更丰富
的细节，在外围区域的动画网格细节逐渐减少。

考虑到过多基础模型的存储需求会导致内存占用的增加，该策略依据注视点位置将可视区域分为三层阶级，
即相应地需要三个基础模型，基于上述提出的时空运动特征驱动的动画网格层次简化算法生成的简化结果进行构建：
高精度层级（LOD0），适用于注视点中心凹区域，确保用户注视区域具有最高的几何保真度和纹理分辨率，满足中央凹视觉对细节辨识的需求。
中精度层级（LOD1），适用于注视点周围区域，在保证基本视觉质量的同时，适当降低几何复杂度，符合人类视觉在中等偏心度下分辨能力衰减的特性。
低精度层级（LOD2）：适用于注视点外围区域，显著降低几何细节和纹理分辨率，匹配人眼在视野外围较低的细节分辨能力，同时最大程度减少计算资源消耗。
各精度层级的简化率参数基于心理物理学实验确定，确保了在不同视觉区域中简化操作不会引起可察觉的视觉质量下降。同时每个细节层级采用统一的数据结构，包含三个核心数据组件：简化后的静态网格几何数据、优化后的动态动画序列数据以及压缩后的
纹理数据。

另外，动态细分层则根据实时注视点位置、视觉感知模型和渲染上下文，在基础LOD模型上执行自适应镶嵌细分。细分程度的控制策略综合考虑两个核心因素：顶点到注视点的偏心距离和时空运动特征重要性，该机制实现了离散LOD层级无法提供的连续细节梯度，精确匹配人眼视觉敏感度的空间分布特性。


% 另外，基于\ref{fig:adaptive_lod}提出的 castleCSF-TS 模型得到的降采样因子，在基础LOD模型上执行自适应镶嵌细分。细分程度的控制策略综合考虑两个核心因素：顶点到注视点的偏心距离和时空运动特征重要性，该机制实现了离散LOD层级无法提供的连续细节梯度，精确匹配人眼视觉敏感度的空间分布特性。


为实现注视点引导的连续细节变化，建立基于视觉感知模型的细分因子计算方法。顶点$v$的基础细分因子$TF_{\text{base}}(v)$根据其到注视点的偏心度$e$计算，采用分段函数形式：
\begin{equation}
TF_{\text{base}}(v) =
\begin{cases}
64, & 0 \leq e \leq e_f \\
32, & e_f < e \leq e_m \\
1, & e > e_m
\end{cases}
\label{eq:base_tessellation_factor}
\end{equation}
其中$e_f$和$e_m$分别为中央凹区域和混合区域的偏心度阈值。为将时空运动特征融入细分决策，基础细分因子需进一步调整：

\begin{equation}
TF_{\text{adjusted}}(v,t) = TF_{\text{base}}(v) \cdot \left[ 1 + \eta \cdot I_(v,t) \right]
\label{eq:adjusted_tessellation_factor}
\end{equation}
其中$\eta$为运动特征调节系数，$I_(v,t)$为顶点$v$在时间$t$的的时空运动特征重要性。该机制确保高运动区域获得更丰富的几何细节，保持动画变形区域的视觉保真度。

在实际渲染过程中，自适应镶嵌细分操作针对每个三角形图元执行，需分别计算边缘细分因子和内部细分因子。对于三角形边$e_{uv}$，其边缘细分因子$TF_{\text{edge}}(e_{uv},t)$通过顶点$u$和$v$的细分因子加权插值得到：
\begin{equation}
TF_{\text{edge}}(e_{uv},t) = w_u \cdot \frac{TF_{\text{adjusted}}(u,t)}{k_u} + w_v \cdot \frac{TF_{\text{adjusted}}(v,t)}{k_v}
\label{eq:edge_tessellation_factor}
\end{equation}

权重系数$w_u$和$w_v$基于顶点到注视点的距离计算：
\begin{equation}
w_u = \frac{d_v}{d_u + d_v}, \quad w_v = \frac{d_u}{d_u + d_v}
\label{eq:distance_weights}
\end{equation}
其中$d_u$和$d_v$分别为顶点$u$和$v$到注视点的三维空间距离，$k_u$和$k_v$为3.2.2 节提出的castle感知模型计算得到的降采样因子，反映顶点所在区域的视觉重要性。三角形内部的细分因子$TF_{\text{inner}}$基于三边长度和最小降采样因子计算：
\begin{equation}
TF_{\text{inner}} = \frac{\max(L_1, L_2, L_3)}{k_{\min}}
\label{eq:inner_tessellation_factor}
\end{equation}
其中$L_1$、$L_2$、$L_3$为三角形三边长度，$k_{\min}$为三个顶点降采样因子的最小值。这种计算方式确保细分后的三角形尺寸与屏幕空间像素分布相匹配，避免过度细分导致的渲染效率损失。

为应对传统LOD方法在切换边界处直接替换几何模型导致明显的视觉跳变问题，本文提出基于过渡区域的混合插值算法。设当前渲染区域涉及从LOD级别$i$向级别$j$的过渡，对于处于过渡区域内的顶点$v$，其最终几何位置通过线性插值计算：
\begin{equation}
p_{\text{final}}(v) = \alpha \cdot p_i(v) + (1-\alpha) \cdot p_j(v)
\label{eq:geometry_interpolation}
\end{equation}

过渡权重$\alpha$基于顶点到LOD切换边界的归一化距离$d_{\text{norm}}$计算：
\begin{equation}
\alpha = \text{smoothstep}\left(\frac{d_{\text{norm}} - d_{\text{start}}}{d_{\text{end}} - d_{\text{start}}}\right)
\label{eq:transition_weight}
\end{equation}
其中$\text{smoothstep}(x) = 3x^2 - 2x^3$为三次平滑函数，$d_{\text{start}}$和$d_{\text{end}}$定义过渡区域的起止边界。这种平滑过渡机制消除了离散LOD切换产生的视觉不连续，在动画播放过程中尤其重要，避免因几何突变导致的运动失真。

对于动画数据的过渡处理，需要同时保证时间维度的连贯性和空间维度的一致性。骨骼权重在LOD过渡期间的插值处理遵循保形原则，确保变形效果的平滑过渡。设顶点$v$在LOD级别$i$和$j$的骨骼权重向量分别为$\mathbf{w}_i$和$\mathbf{w}_j$，过渡期间的权重向量计算为：
\begin{equation}
w_{\text{transition}}(v) = \alpha \cdot w_i(v) + (1-\alpha) \cdot w_j(v)
\label{eq:weight_interpolation}
\end{equation}

插值后的权重向量需进行归一化处理以保证$\sum_k w_k = 1$。动画关键帧的过渡处理采用时间对齐的插值方法，确保不同LOD级别的动画序列在过渡期间保持同步。对于时间$t$处的姿势插值，首先将相邻LOD级别的动画序列对齐到相同时间基准，然后基于过渡权重进行混合：
\begin{equation}
P_{\text{transition}}(t) = \alpha \cdot P_i(\phi_i(t)) + (1-\alpha) \cdot P_j(\phi_j(t))
\label{eq:pose_interpolation}
\end{equation}

其中$P_i$和$P_j$分别为LOD级别$i$和$j$的骨骼姿势函数，$\phi_i(t)$和$\phi_j(t)$为时间映射函数，确保不同简化率的动画序列在相同时间点具有对应姿势。










% \subsection{感知引导的动画网格简化算法}

% 注视点渲染技术需要解决的一个关键问题是如何用不同的质量渲染不同的区域。
% 对于该问题，计算机图形学中的 LOD 技术为绘制由包含不同数据量的几何网格组成的三维场景提供了一种解决方案。
% 该技术根据相机与模型对象之间的距离动态调整对象的细节级别，通常，离相机较远的对象使用较低的细节级别，而近处的对象则使用较高的细节级别。
% 根据不同的实现方式和应用场景，传统的 LOD技术可分为离散 LOD 和连续 LOD。

% 离散 LOD 算法的基本原理是预先为对象生成多个不同精度的模型，并根据视物距离选择相应级别的模型渲染，简单易实现。但该方法存在一个严重的问题，即当距离变化导致层次切换时，可能出现明显的视觉不连贯。
% 连续 LOD 算法，也被视为一种动态 LOD，可以根据视物距离实时生成或修改几何细节，以确保视觉效果的平滑过渡。
% 然而，对于复杂场景的渲染而言，实时计算细节变化可能会引入延迟，降低用户体验感；同时，在原始精度模型上不断简化会出现模型失真问题。

% 因此，本方法对传统的 LOD 技术进行改进，将视物距离和降采样因子共同作为关键指标动态调整模型细节。
% 引入镶嵌技术，并基于感知模型计算得到的降采样因子动态调整细分程度，在不同区域执行不同程度的曲面细分，提出注视点驱动的镶嵌技术，进而提出感知自适应细分 LOD 算法，该算法框架如图 3.4 所示。


% ASLoD 结合了离散 LOD 技术和镶嵌技术。首先，使用离散 LOD 根据视物距离初始分阶处理网格模型；
% 其次，在对应 LOD 层次内使用镶嵌技术对注视区域的模型执行网格细分，降采样因子动态调整细分程度，使得在中央凹区域的模型拥有更丰富的细节，在外围区域的模型细节逐渐减少。
% 镶嵌技术被用来调整网格的密度，将原始低多边形网格细分为更精细的网格，相较于基于点的方法细分效果更好。
% 该技术依赖于 GPU 的镶嵌单元，可用于实时性能优化。

% \subsection{时空感知重要性度量模型}

% 为量化动画网格中不同区域的视觉重要性，本文提出时空感知重要性度量模型。该模型综合考虑了空间视觉特征、时间运动特征和结构关联特征，定义顶点$v$在时间$t$的感知重要性$I(v,t)$为：

% \begin{equation}
%     I(v,t) = \omega_c \cdot C(v,t) + \omega_m \cdot M(v,t) + \omega_g \cdot G(v) + \omega_j \cdot J(v)
%     \label{eq:perceptual_importance}
% \end{equation}

% 其中各分量定义如下：对比敏感度分量 $ C(v,t) $ 基于第三章提出的 castleCSF-TS 模型计算，用于衡量顶点在当前视觉上下文中的显著程度，
% 该模型综合考虑了空间频率、时间频率和偏心度等因素。

% 运动显著性分量 $ M(v,t) $ 表征顶点在时间维度上的运动强度，通过计算时间窗口内的平均速度来度量，
% 其推导基于顶点位置序列的差分：
% \begin{equation}
% M(v,t) = \frac{1}{T} \sum_{\tau=t-T}^{t} \frac{\| v(\tau) - v(\tau-1) \|}{\Delta \tau}
% \label{eq:motion_saliency}
% \end{equation}
% 其中，$ T $ 为时间窗口大小，本文取 $ T=5 $ 帧，以确保对短期运动模式的捕捉；$ \Delta \tau $ 为时间间隔，
% 通常取帧率倒数，该分量强调了快速运动的顶点更易吸引视觉注意。

% 几何重要性分量 $ G(v) $ 基于局部曲率和法线变化来度量静态几何特征，反映顶点在形状结构中的关键性，
% 其定义为：
% \begin{equation}
% G(v) = \alpha \cdot \kappa(v) + \beta \cdot \frac{1}{|N(v)|} \sum_{u \in N(v)} (1 - \vec{n}_v \cdot \vec{n}_u)
% \label{eq:geometric_importance}
% \end{equation}
% 式中，$ \kappa(v) $ 为顶点曲率，描述局部弯曲程度；$ N(v) $ 为邻接顶点集合；$ \vec{n}_v $ 和 $ \vec{n}_u $ 分别为顶点 $ v $ 及其邻点 $ u $ 的法向量；
% $ \alpha $ 和 $ \beta $ 为权重参数，用于平衡曲率和法线变化的贡献，高曲率或法线突变的区域往往对应几何细节丰富的区域。

% 关节关联度分量 $ J(v) $ 反映顶点与骨骼关节的关联强度，用于动画或角色模型中重点关注关节附近区域，
% 该分量通过高斯函数计算：
% \begin{equation}
% J(v) = \exp\left(-\frac{\min_{j \in \mathcal{J}} \|v - j\|^2}{2\sigma^2}\right)
% \label{eq:joint_association}
% \end{equation}
% 其中，$ \mathcal{J} $ 为关节集合，$ \sigma $ 为尺度参数，控制关联强度的衰减速率，该公式表明顶点越接近最近关节，
% 其关联度越高，视觉重要性也相应提升。

% 权重系数$\omega_c, \omega_m, \omega_g, \omega_j$通过心理物理学实验标定，取值分别为0.40、0.30、0.20、0.10。


% 基于上述时空感知重要性度量，本文提出多特征协同简化策略。传统网格简化算法如二次误差度量（QEM）主要基于几何误差进行边折叠决策，往往忽略了法线、纹理等视觉关键特征的保持效果。针对这一局限性，本研究提出一种基于平均偏差的多特征加权简化算法，通过在简化过程中综合考虑位置坐标、法线向量和纹理坐标等多个特征维度，实现更好的视觉保真度。

% 在平均偏差计算模型中，对于任意相邻顶点 $V_i$ 和 $V_j$，定义其多特征平均偏差 $D_{ij}$ 的计算公式如下：
% \begin{equation}
%     D_{ij} = \frac{1}{n} \sum_{k=1}^{n} w_k \cdot d_k(F_i^k, F_j^k)
%     \label{eq:average_deviation}
% \end{equation}
% 其中，$n$ 表示特征维度总数，$F_i^k$ 和 $F_j^k$ 分别代表顶点 $V_i$ 和 $V_j$ 在第 $k$ 个特征维度上的取值，$d_k(\cdot, \cdot)$ 为第 $k$ 个特征的距离度量函数，$w_k$ 为对应的特征权重系数。具体而言，本文考虑三个关键特征维度：位置坐标（权重系数 $w_{\text{pos}}$）、法线向量（权重系数 $w_{\text{norm}}$）和纹理坐标（权重系数 $w_{\text{tex}}$）。

% 各特征的距离度量函数分别定义为：位置特征采用欧氏距离度量，该度量反映了顶点在三维空间中的实际几何距离，计算公式为：
% \begin{equation}
%     d_{\text{pos}} = \sqrt{(x_i - x_j)^2 + (y_i - y_j)^2 + (z_i - z_j)^2}
%     \label{eq:position_distance}
% \end{equation}
% 其中，$(x_i, y_i, z_i)$ 和 $(x_j, y_j, z_j)$ 分别表示顶点 $V_i$ 和 $V_j$ 在三维笛卡尔坐标系中的坐标值。该距离度量能够准确刻画两个顶点在三维空间中的实际分离程度，是衡量几何相似性的基础指标。

% 法线特征考虑单位向量之间的夹角因素，采用点积运算进行度量，其计算公式为：
% \begin{equation}
%     d_{\text{norm}} = 1 - \vec{n}_i \cdot \vec{n}_j
%     \label{eq:normal_distance}
% \end{equation}
% 其中，$\vec{n}_i$ 和 $\vec{n}_j$ 分别表示顶点 $V_i$ 和 $V_j$ 处的单位法线向量。由于法线向量已归一化，其点积结果等于两法线向量夹角的余弦值，即 $\vec{n}_i \cdot \vec{n}_j = \cos\theta$，其中 $\theta$ 为两法线向量之间的夹角。当两法线方向完全一致时，点积值为 $1$，差异度量为 $0$；当两法线相互垂直时，点积值为 $0$，差异度量为 $1$；当两法线方向相反时，点积值为 $-1$，差异度量为 $2$。该度量能够有效反映网格表面曲率变化的连续性。

% 纹理特征则采用曼哈顿距离度量，计算公式为：
% \begin{equation}
%     d_{\text{tex}} = |u_i - u_j| + |v_i - v_j|
%     \label{eq:texture_distance}
% \end{equation}
% 其中，$(u_i, v_i)$ 和 $(u_j, v_j)$ 分别表示顶点 $V_i$ 和 $V_j$ 在纹理空间中的 UV 坐标。纹理坐标通常规范化到 $[0,1]$ 区间，表示顶点在纹理图像上的映射位置。曼哈顿距离通过计算两个顶点在纹理空间 $U$ 方向和 $V$ 方向上的坐标差异绝对值之和，提供了纹理对齐程度的量化指标，该度量对于保持纹理连续性和避免纹理扭曲具有重要作用。

% 在边折叠策略的设计中，算法优先处理平均偏差值较小的边，这一机制确保视觉特征相似的区域能够优先被简化。同时，引入基于视觉重要性的权重调整机制，对角色关节区域、面部五官等高视觉关注度区域赋予更高的保留优先级。具体实现中，通过计算每个顶点的局部曲率权重和关节关联度，建立视觉重要性评估函数，动态调整各顶点在简化过程中的被折叠概率，从而在整体简化过程中有效保持模型的视觉关键特征。

% 在纹理数据优化方面，本研究将网格几何简化过程与纹理优化紧密结合，采用 Mipmap 金字塔结构和自适应可伸缩纹理压缩（ASTC）技术。针对不同细节层级的网格模型，分配合适的纹理分辨率：高细节层级使用原始分辨率或轻度压缩的纹理格式，以保持视觉细节；低细节层级则采用高强度压缩格式，显著减少显存占用。这种多分辨率纹理管理方案在保证视觉质量的前提下，实现了纹理存储效率的显著提升。

% 本算法通过综合考虑几何特征与纹理特征的多维相似性，建立了一套完整的网格简化评估体系。与传统方法相比，不仅关注几何误差的控制，更注重视觉关键特征的保持，从而在相同简化率条件下获得更好的视觉保真度。实验结果表明，该算法在复杂角色模型的简化过程中，能够有效保持关节部位和面部区域的细节特征，为后续的实时渲染应用提供了高质量的简化模型。

% 在虚拟现实系统中，动态动画序列的高效处理是保证实时渲染性能的关键技术挑战。传统动画序列优化方法往往侧重于单一维度的压缩，难以在保持视觉质量的同时实现高压缩率。为此，本文提出一种基于运动连贯性分析的双阶段动画序列优化算法，通过关键帧压缩和骨骼数据精简两个阶段的协同优化，在严格保持动画视觉质量的前提下，显著降低动态数据的存储需求和计算负载。

% 第一阶段基于骨骼运动轨迹分析实现关键帧压缩，该过程建立在对骨骼运动轨迹连续性的分析基础上。算法通过量化相邻关键帧间的骨骼位移差异，识别并剔除视觉冗余的关键帧，同时保证重要运动特征的完整保留。对于骨骼系统中特定骨骼 $B_i$ 在关键帧 $t_j$ 至 $t_{j+1}$ 时间区间内的运动轨迹，定义其平均位移误差度量函数为：
% \begin{equation}
%     E_{ij} = \frac{1}{N} \sum_{k=1}^{N} \left\| T_{ij}(p_k) - T_{i(j+1)}(p_k) \right\|
%     \label{eq:average_displacement_error}
% \end{equation}
% 其中，$p_k$ 代表受骨骼影响的顶点集合，$T_{ij}$ 表示骨骼 $B_i$ 在关键帧 $t_j$ 的空间变换矩阵，$N$ 为该骨骼影响的顶点总数。该误差项有效表征了相邻关键帧间骨骼运动的视觉显著程度。

% 为确定能够保证视觉无损的位移误差阈值 $\theta_{\text{key}}$，本研究设计了系统的心理物理学实验。实验采用双盲对照设计，招募20名矫正视力正常（$\geq1.0$）的健康成年人（10男10女），年龄20-35岁，均无视觉感知障碍史。实验设备采用HTC Vive Pro Eye头显（144Hz刷新率，110°视场角），配合Tobii Pro Nano眼动仪（250Hz采样率，0.5°精度）记录眼动数据。实验刺激为一段标准化的虚拟角色动画序列（包含行走、奔跑、跳跃等典型动作），时长为30秒，原始关键帧率为60 FPS。

% 实验采用自适应阶梯法确定每位参与者的感知阈值。具体流程为：每次试验中，向参与者随机呈现原始动画序列与经过不同阈值压缩的动画序列（压缩阈值从1.0mm开始），要求参与者判断两段动画是否存在视觉差异。根据参与者的判断结果动态调整压缩阈值：若参与者报告“无差异”，则在下一试次中增加压缩强度（即增大阈值）；若报告“有差异”，则降低压缩强度。初始步长为0.2mm，经过4次反转后步长减小至0.05mm。每位参与者完成8次反转后结束测试，取最后6次反转阈值的算术平均值作为该参与者的个体感知阈值。

% 实验结果显示，20名参与者的平均感知阈值为0.48mm。基于此，为涵盖95%置信区间并确保算法鲁棒性，确立位移误差阈值为 $\theta_{\text{key}} = 0.5$mm。当 $E_{ij} < \theta_{\text{key}}$ 时，判定关键帧 $t_{j+1}$ 在视觉感知上存在冗余，可予以安全剔除。这一严格阈值的确立确保了压缩过程仅移除那些不会引起视觉感知差异的关键帧，从根本上保证了动画质量的视觉保真度。

% 为保证压缩后动画序列的连续性和运动平滑性，对保留的关键帧采用Catmull-Rom样条插值算法生成中间过渡帧。插值过程数学表述如下：
% \begin{equation}
%     P(t) = P_1 + t^2 \left[ -P_0 + P_2 + t\left( 2P_0 - 5P_1 + 4P_2 - P_3 + t\left( -P_0 + 3P_1 - 3P_2 + P_3 \right) \right) \right]
%     \label{eq:catmull_rom_interpolation}
% \end{equation}
% 其中，$P_0$、$P_1$、$P_2$、$P_3$ 构成四个连续的控制点序列，$t \in [0,1]$ 为归一化的时间参数。同时该插值方法具有$C^1$连续性特性，能够生成符合自然运动规律的平滑轨迹，有效维持原始动画的视觉真实感。

% 第二阶段基于权重分布相似性聚类实现骨骼数据精简，该阶段通过分析骨骼权重分布的统计相似性特征，识别并合并功能冗余的骨骼节点。定义骨骼 $B_i$ 和 $B_j$ 之间的权重分布相似性度量函数为：
% \begin{equation}
%     \text{Sim}(B_i, B_j) = \frac{W_i \cdot W_j}{\|W_i\| \|W_j\|}
%     \label{eq:weight_similarity}
% \end{equation}
% 其中，$W_i$ 和 $W_j$ 分别表示骨骼 $B_i$ 和 $B_j$ 的权重向量，该向量完整描述了各骨骼对网格顶点的变形影响程度。

% 为确定合适的相似度阈值 $\theta_{\text{bone}}$，本研究进行了第二项心理物理学实验。实验采用相同的参与者群体和设备配置。实验刺激为经过不同相似度阈值（从0.80到0.99，步长0.01）精简骨骼系统后的动画序列，每个序列持续20秒，包含复杂的多关节协同动作。实验采用二项迫选法，每次试验中同时呈现原始动画和精简后动画（顺序随机），要求参与者判断哪一段动画看起来更自然、无异常变形。实验采用恒刺激法，每个阈值水平呈现10次，共200次试验（分4组完成，组间休息10分钟）。

% 实验结果表明，当相似度阈值 $\theta_{\text{bone}} \geq 0.90$ 时，参与者已难以可靠区分原始与精简动画。基于此，将相似度阈值设定为 $\theta_{\text{bone}} = 0.90$。当 $\text{Sim}(B_i, B_j) > \theta_{\text{bone}}$ 时，判定两个骨骼在变形功能上具有高度相似性，可将其合并为单个骨骼节点。合并过程中，采用加权平均算法重新计算新骨骼的权重分布：
% \begin{equation}
%     W_{\text{new}} = \frac{\|W_i\| \cdot W_i + \|W_j\| \cdot W_j}{\|W_i\| + \|W_j\|}
%     \label{eq:weight_fusion}
% \end{equation}
% 其中，$W_i$ 和 $W_j$ 分别表示两个骨骼的权重向量，$\|W_i\|$ 和 $\|W_j\|$ 表示对应权重向量的范数，$W_{\text{new}}$ 表示合并后的新权重向量。该权重融合策略确保了骨骼合并后网格变形效果的视觉一致性，有效避免了因骨骼精简可能引起的模型变形异常。

% 本算法相较于传统动画序列优化方法具有显著优势：首先，通过双阶段优化架构分别处理时序冗余和结构冗余，实现了更全面的数据压缩效果；其次，基于运动连贯性的误差度量方法结合精确的视觉感知阈值，确保了优化过程中重要运动特征的完整保留；最后，严格的数学建模和参数验证流程，保证了算法在实际应用中的可靠性和稳定性。



% \subsection{感知自适应细分算法}

% 感知自适应细分算法的设计核心在于将人类视觉系统的感知特性与实时渲染的效能需求相结合，构建一种能够动态调整网格细分程度的智能化机制。该算法以castleCSF-TS视觉感知模型为理论基础，通过量化不同视觉区域的细节敏感度，为每个顶点计算自适应的细分因子，从而在注视点区域提供高精度细节，在外围区域则通过降低细分程度来节省计算资源，同时确保不引起视觉感知上的质量下降。算法的输入包括当前帧的动画网格数据、注视点位置、实时性能预算以及预计算的感知重要性图，输出则为每个顶点的细分因子，这些因子将直接指导GPU细分管线进行自适应的网格细分。

% 算法的数学基础建立在对人类视觉系统空间非均匀特性的深刻理解之上。视网膜的视觉敏感度分布呈现以中央凹为中心向周边指数衰减的趋势，这一特性可通过偏心度$e$来量化，即顶点在屏幕空间中的投影位置与注视点之间的角度偏差。基于这一生理特性，我们将渲染区域划分为三个层次：中央凹区域（$e \leq 2^\circ$）、混合区域（$2^\circ < e \leq 10^\circ$）和外围区域（$e > 10^\circ$）。每个区域对应不同的基础细分因子，形成分段函数：

% \begin{equation}
% TF_{\text{base}}(e) = 
% \begin{cases} 
% 64, & 0 \leq e \leq e_f \\
% 32, & e_f < e \leq e_m \\
% 1, & e > e_m 
% \end{cases}
% \label{eq:base_tess_factor}
% \end{equation}

% 其中$e_f = 2^\circ$和$e_m = 10^\circ$为区域边界阈值。这一划分确保了中央凹区域获得最高细分精度，以满足该区域对细节辨识的高要求；混合区域采用中等细分精度，在视觉质量与计算效率之间取得平衡；外围区域则利用人类视觉系统的感知冗余，将细分程度降至最低。

% 然而，仅基于偏心度的细分策略未能充分考虑顶点自身的视觉重要性差异。因此，我们引入感知重要性调整机制。顶点$v$的感知重要性$I(v)$由公式(\ref{eq:perceptual_importance})计算得到，该公式融合了对比敏感度、运动显著性、几何重要性和关节关联度四个维度的信息。通过线性插值将感知重要性映射到细分因子的调整系数上，得到调整后的细分因子：

% \begin{equation}
% TF_{\text{adj}}(v) = TF_{\text{base}}(e_v) \cdot \left[1 + \alpha \cdot \frac{I(v) - I_{\min}}{I_{\max} - I_{\min}} \right]
% \label{eq:adjusted_tess_factor}
% \end{equation}

% 其中$\alpha = 0.5$为调整系数，用于控制感知重要性对细分因子的影响强度。这一步骤确保了视觉上更重要的顶点（如快速运动的关节区域或高曲率表面）能够获得更高的细分因子，从而在简化过程中得到更好的保护。

% 考虑到动画网格的动态特性，顶点的运动状态同样影响细分需求。我们引入运动敏感度补偿因子，以增强高运动区域的细分程度。运动补偿后的细分因子计算如下：

% \begin{equation}
% TF_{\text{motion}}(v) = TF_{\text{adj}}(v) \cdot \left[1 + \beta \cdot M(v,t) \right]
% \label{eq:motion_adjusted_tess_factor}
% \end{equation}

% 其中$M(v,t)$为顶点$v$在时间$t$的运动显著性，由公式(\ref{eq:motion_saliency})定义；$\beta = 0.3$为运动补偿系数。这一补偿机制确保在动画播放过程中，高变形区域能够获得足够的细分精度，避免因简化而产生的运动失真。

% 为满足实时渲染的性能要求，算法还需根据当前帧率进行动态调整。我们引入基于帧率反馈的性能约束机制：若当前帧率$F_{\text{current}}$低于目标帧率$F_{\text{target}}$（通常设为90 FPS），则按比例降低细分因子；否则保持原细分因子不变。具体调整公式为：

% \begin{equation}
% TF_{\text{final}}(v) = 
% \begin{cases} 
% TF_{\text{motion}}(v), & F_{\text{current}} \geq F_{\text{target}} \\
% TF_{\text{motion}}(v) \cdot \left( \frac{F_{\text{current}}}{F_{\text{target}}} \right)^{\gamma}, & F_{\text{current}} < F_{\text{target}}
% \end{cases}
% \label{eq:final_tess_factor}
% \end{equation}

% 其中$\gamma = 2.0$为性能衰减指数。这一机制保证了在性能紧张时，算法能够自适应地降低细分负载，从而维持流畅的交互体验。

% 在得到每个顶点的最终细分因子后，算法进一步计算每条边的细分因子。对于连接顶点$u$和$v$的边$e_{uv}$，其细分因子通过两端点细分因子的加权平均得到：

% \begin{equation}
% TF_{\text{edge}}(e_{uv}) = \omega_u \cdot \frac{TF_{\text{final}}(u)}{k_u} + \omega_v \cdot \frac{TF_{\text{final}}(v)}{k_v}
% \label{eq:edge_tess_factor}
% \end{equation}

% 权重$\omega_u$和$\omega_v$基于顶点到注视点的距离计算，使得距离注视点更远的顶点对边细分因子的影响更小：

% \begin{equation}
% \omega_u = \frac{d_v}{d_u + d_v}, \quad \omega_v = \frac{d_u}{d_u + d_v}
% \label{eq:weight_computation}
% \end{equation}

% 其中$d_u$和$d_v$分别是顶点$u$和$v$到注视点的三维空间距离，$k_u$和$k_v$是顶点的降采样因子，反映了在屏幕空间中的采样需求。这一设计确保了细分因子在空间上的平滑过渡，避免了因突变而产生的视觉瑕疵。

% 对于三角形面片的内部细分因子，我们采用基于最大边长和最小降采样因子的计算方法：

% \begin{equation}
% TF_{\text{inner}}(t) = \frac{L_{\max}}{k_{\min}}
% \label{eq:inner_tess_factor}
% \end{equation}

% 其中$L_{\max}$是三角形三边中的最大长度，$k_{\min}$是三个顶点降采样因子的最小值。这一公式确保在面片内部，细分程度能够适应几何特征的变化，在保持视觉连续性的同时避免不必要的细分。

% 在细分执行阶段，算法将计算得到的细分因子传递给GPU的细分着色器。现代图形渲染管线中的细分着色器包括外壳着色器、细分器和域着色器。在外壳着色器中，我们为每个面片设置边细分因子和内部细分因子；细分器根据这些因子生成新的几何图元；域着色器则对每个新顶点进行插值计算，包括位置、法线、纹理坐标以及骨骼权重。对于动画网格，骨骼权重的正确插值至关重要。新顶点$v_{\text{new}}$的骨骼权重通过原始三角形顶点的权重重心插值得到：

% \begin{equation}
% w_{\text{new}}^j = \sum_{i=1}^{3} \lambda_i \cdot w_i^j, \quad j = 1, \ldots, B
% \label{eq:bone_weight_interpolation}
% \end{equation}

% 其中$\lambda_i$是重心坐标分量，$w_i^j$是原始顶点$i$对骨骼$j$的权重。插值后需进行归一化处理以确保权重之和为1：

% \begin{equation}
% \tilde{w}_{\text{new}}^j = \frac{w_{\text{new}}^j}{\sum_{k=1}^{B} w_{\text{new}}^k}
% \label{eq:weight_normalization}
% \end{equation}

% 为进一步优化，我们移除权重值小于阈值$\epsilon = 0.01$的骨骼影响，并将这些权重重新分配给主要骨骼，从而减少骨骼数量并保持变形效果的视觉一致性。



% \section{实验设计与结果分析}
% \subsection{实验环境}


% 为验证本文提出的感知自适应动态LOD技术的可行性与有效性，搭建了一套集成高精度眼动追踪与感知模型驱动的虚拟现实渲染实验平台。该平台由硬件系统与软件框架组成，具体配置如表\ref{tab:experiment-configuration}所示。

% \begin{table}[htbp]
%   \centering
%   \caption{实验软硬件环境配置}
%   \label{tab:experiment-configuration}
%   \renewcommand{\arraystretch}{1.3}
%   \begin{tabularx}{\textwidth}{>{\centering\arraybackslash}m{1.5cm} >{\raggedright\arraybackslash}X >{\raggedright\arraybackslash}X}
%     \toprule
%     \textbf{类别} & \textbf{名称} & \textbf{配置/版本} \\
%     \midrule
%     \textbf{硬件} & 中央处理器 (CPU) & Intel Core i9-13900K \\
%     & 图形处理器 (GPU) & NVIDIA GeForce RTX 4090 \\
%     & 内存 & 64GB DDR5 (2×32GB) \\
%     & 主VR头显 & HTC Vive Pro Eye  \\
%     & 眼动追踪模块 & 内置Tobii眼动仪 \\
%     & 辅助高精度眼动仪 & SR Research EyeLink 1000 Plus \\
%     & 定位基站 & SteamVR Base Station 2.0 \\
%     \midrule
%     \textbf{软件} & 操作系统 & Windows 11 Professional \\
%     & 渲染与开发引擎 & Unity 2022.3 LTS \\
%     & 眼动追踪SDK & SRanipal SDK v1.3.2 \\
%     & 数据分析 & MATLAB R2024a \\
%     \bottomrule
%   \end{tabularx}
% \end{table}


% \label{sec:experiment-environment}

% \subsection{评价指标}
% 本方法改进了传统 LOD 算法，首先基于原始模型简化出八个不同精度的基础模
% 型，根据视物距离动态选定对应级别的基础模型，其次根据注视点的位置动态地对基
% 础模型不同区域进行不同程度的细分。由于基于传统注视点渲染设置发生 LOD 切换
% 会产生几何突变、纹理闪烁以及时序不一致性等问题，导致渲染质量损失可感知，因
% 此，在本实验中使用每帧 LOD 切换次数这一指标来衡量本算法应用于注视点渲染时
% 的质量可感知损失改善，该指标是指一帧时间内场景中所有物体发生 LOD 层级切换
% 的总次数。每帧 LOD 切换次数的数学公式可表示为：
% \begin{equation}
% C_{\text{switch}}(t) = \sum_{i=1}^n I(\text{LOD}_i(t) \neq \text{LOD}_i(t-1))
% \tag{3-34}
% \end{equation}

% 其中，\( C_{\text{switch}}(t) \) 表示第 \( t \) 帧的 LOD 切换次数，\( n \) 表示场景中物体总数，\( \text{LOD}_i(t) \) 表示物体在第 \( i \) 帧的细节层级，\( I(\cdot) \) 表示指示函数，条件成立时值为 1，否则为 0。

% 当每帧 LOD 切换次数显著降低，表明系统优化了细节层级的动态调整策略。具体而言，LOD 切换频率的降低延长了同一细节层级的持续时间，使模型简化与细化的过渡节奏适配人眼视觉暂留效应。在此阈值内，连续帧间的画面变化可被视觉系统融合为平滑运动；然而，若切换频率过高，则可能突破融合阈值，导致几何突变与纹理闪烁等伪影被明显感知。

% 在选定 LOD 层级后，对于注视点所在的中央凹区域执行较大程度的细分从而满足人眼视觉感知，而对于偏心度逐渐增大的混合区域以及外围区域则只需较小程度的细分，或者不需要细分，细分的程度主要由感知模型输出的最大降采样因子决定。有效细分触发率是衡量感知驱动 LOD 算法决策准确性的关键指标，用于评估算法在动态注视点渲染中，根据感知模型正确触发细分操作的能力。该指标的数学定义为：

% \begin{equation}
% \delta_{AT} = \frac{1}{N} \sum_{n=1}^N I(k_n \geq k(e)) \times 100\%
% \tag{3-35}
% \end{equation}

% 其中，\( N \) 表示场景中需要处理的网格模型总数；\( k_n \) 表示模型 \( n \) 的实际细分因子，\( k(e) \) 表示由感知模型计算得到的理论细分阈值，且该值与模型所处偏心度有关；\( I(\cdot) \) 为指示函数，当条件满足时取 1，否则取 0。

% 通过对比实际细分因子与感知模型的理论阈值，量化算法在动态渲染中响应视觉需求的能力。高触发率，即 \( \delta_{AT} \geq 98\% \)，表明算法能准确识别注视点区域的高细节需求，确保关键区域的视觉保真度。同时，更高的触发率意味着能够更加平滑地过渡混合区域以减少视觉不连贯性，并更有效地利用计算资源以提高渲染效率。

% 冗余三角面压缩率是评估注视点渲染中 LOD 算法模型简化效率的核心指标。指标通过量化算法对不同视物距离下的注视点区域冗余几何数据的压缩能力，直接反映了算法在降低渲染负载方面的有效性。其数学定义为：

% \begin{equation}
% \delta_{LOD} = \left( 1 - \frac{N_{\text{output}}}{N_{\text{input}}} \right) \times 100\%
% \tag{3-36}
% \end{equation}

% 其中，\( N_{\text{input}} \) 表示原始模型的三角面数，\( N_{\text{output}} \) 表示算法处理后的实际渲染模型三角面数。压缩率越高，表明算法对模型的简化程度越大，能显著减少 GPU 处理的几何数据量，从而加速渲染。

% 另外，渲染时间稳定性是衡量注视点渲染系统实时性能的核心指标，其不仅反映了算法对硬件资源的利用效率，更直接决定虚拟现实体验的流畅性。VR 系统要求帧
% 率稳定在 90FPS 以上，否则引发视觉卡顿与晕动症。

% \subsection{实验设计与结果分析}
% 本实验在 Unity 2021.3 开发环境中，选用包含动态植被、地形位移贴图和动态天
% 气系统的开放地形作为测试场景。呈现该复杂地形的外貌需要庞大的地形网格，在不
% 使用 LOD 技术时，应用阶段需提供高细节网格模型，如图 3.7 所示。而使用 LOD 技
% 术并结合镶嵌技术，应用阶段只需提供对应阶级简化的地形网格轮廓，细分阶段动态添加网格顶点，将应用阶段 CPU 提供的低细节网格转换为几何阶段 GPU 上感知自适
% 应的高细节网格，从而显著减少顶点着色器处理的顶点数量，如图 3.8 所示。

% 本实验为对比实验，将本文提出的注视点渲染算法与静态设置的注视点渲染算法
% 进行比较。本文算法基于感知模型动态划分渲染区域并执行相应程度的曲面细分，而
% 静态设置的算法则使用传统的区域划分方式对注视点区域执行特定程度的曲面细分。
% 静态设置实验组通过自定义脚本实现固定阈值分层，使用线性比例代表视锥
% 体衰减，将渲染区域按固定偏心度划分为中央凹区域（注视点±3°）、近中央凹区域
% （3°-5°）和外围区域（＞5°，至视口边界），并基于最精简的基础模型在各区域执行
% 不同程度的曲面细分。本文算法实验组则使用基于感知模型的自适应细分 LOD 方法，
% 首先在 ComputeShader 中实时计算不同条件下基于 BD-castleCSF 模型的对比敏感度
% 阈值，并根据该阈值动态划分中央凹区域和混合区域的大小，同时计算得到动态细分
% 因子在基础模型上执行曲面细分。
% 本实验选取视野中左上角和中间偏下的区域作为注视区域，以便测试算法在不同
% 网格密度下的渲染结果及性能，两种方法渲染结果分别如图 3.9 和图 3.10 所示。


% 根据图 3.9(b)和图 3.9(d)可知，静态设置细分因子的方法仅适用于屏幕上特定尺
% 寸的渲染图元，当位于中央凹区域的图元大于该尺寸时，细分不足；反之则导致过度
% 细分。与静态设置算法相比，本文算法综合了考虑刺激维度对视觉的影响，通过计算
% 对比敏感度动态划分渲染区域并调整细分因子，更符合注视点渲染要求，实现了动态
% 渲染。两组实验的渲染性能统计结果如表 3.4 所示。两组实验结果表明，相较于传统
% 的静态设置方案，感知自适应的方法在域着色器顶点优化方面展现出显著优势，降幅
% 达 58\%，但因动态评估引入实时计算负载，渲染耗时平均增加约 12.2\%。尽管固定
% 区域划分方法渲染速度稍快，但会导致资源分配不合理，浪费计算资源。


% 再者，由表 3.5 可知，本文算法的 LOD 切换频率降低了 67\%，显著减少了几何
% 突变和纹理闪烁等视觉伪影，有效改善了渲染质量的可感知损失，降低了 GPU 频繁
% 切换状态而产生的计算开销。使用本文算法有效细分触发率达 98.7\%，提升了 19.9\%，
% 更精准地在需要平滑过渡的区域执行曲面细分，体现了该算法对感知需求的精准响应。
% 而使用静态设置算法固定半径与动态场景不匹配，导致部分外围高对比度物体未被正
% 确细分。对于三角面压缩率，相比静态设置算法，本文算法通过动态调整阈值，避免
% 了固定分层对中距离物体的过度保留。本文算法实验组的帧时间标准差较静态设置算
% 法实验组有所降低，表明渲染负载的时序分布更加均匀，证明了视觉连贯性的提升。


% 上述实验结果表明，本文提出的基于感知模型的自适应细分 LOD 方法将注视点
% 渲染的分层策略从经验驱动的固定阈值转向感知模型驱动的动态变化，在减少细节层
% 次频繁切换的同时，精准提升了关键区域的渲染精度，有效抑制了视觉伪影和几何突
% 变。此外，本方法的感知自适应的资源调度机制兼顾显存与算力平衡，既避免了过度
% 简化导致的质量可感知损失，又消除了冗余计算引发的帧率波动，最终在视觉保真度
% 与渲染效率之间实现了更优平衡。

% \begin{table}[htbp]
%     \centering
%     \caption{注视点驱动的渲染性能统计结果}
%     \label{tab:gaze_rendering_performance}
%     \renewcommand{\arraystretch}{1.5}
%     \begin{tabularx}{\textwidth}{
%         >{\centering\arraybackslash}m{0.18\textwidth}
%         >{\centering\arraybackslash}m{0.15\textwidth}
%         >{\centering\arraybackslash}m{0.16\textwidth}
%         >{\centering\arraybackslash}m{0.16\textwidth}
%         >{\centering\arraybackslash}m{0.15\textwidth}}
%         \toprule
%         \textbf{渲染方法} & \textbf{注视点区域} & \textbf{顶点着色器接收顶点数} & \textbf{域着色器输出顶点数} & \textbf{渲染时间（单位：微秒）} \\
%         \midrule
%         静态设置算法 & 左上角 & 1,279,126 & 685,145 & 6325.33 \\
%                      & 中下角 & 1,279,126 & 550,024 & 5012.70 \\
%         \midrule
%         \multirow{2}{*}{本文算法实验组} & 左上角 & 1,279,126 & 613,980 & 6623.74 \\
%                      & 中下角 & 1,279,126 & 537,232 & 5974.61 \\
%         \bottomrule
%     \end{tabularx}
% \end{table}

% \begin{table}[htbp]
%     \centering
%     \caption{评价指标分析}
%     \label{tab:evaluation_metrics}
%     \renewcommand{\arraystretch}{1.5}
%     \begin{tabularx}{\textwidth}{
%         >{\centering\arraybackslash}m{0.4\textwidth}
%         >{\centering\arraybackslash}m{0.3\textwidth}
%         >{\centering\arraybackslash}m{0.3\textwidth}}
%         \toprule
%         \textbf{评价指标} & \textbf{静态设置算法实验组} & \textbf{本文算法实验组} \\
%         \midrule
%         平均每帧 LOD 切换次数 \( C_{\text{switch}} \) & 9.8 & \textbf{3.2} \\
%         有效细分触发率 \( \delta_{AT} \) & 82.3\% & \textbf{98.7\%} \\
%         冗余三角面压缩率 \( \delta_{LOD} \) & 57.1\% & \textbf{63.9\%} \\
%         帧时间标准差 \( \delta_t \) & 3.2 & \textbf{2.1} \\
%         \bottomrule
%     \end{tabularx}
% \end{table}



































